\section{Setup}\label{sec:setup}

% hand-written study overview (no new numbers; every value appears in the generated tables or result READMEs)
\begin{table}[t]
\centering
\small
\caption{Study overview. GPT-2 and Gemma use different SAE families, hook conventions and data, so their results are
reported side by side and never ranked against each other.}
\label{tab:overview}
\begin{adjustbox}{max width=\linewidth}
\begin{tabular}{lll}
\toprule
 & GPT-2 Small & Gemma-2-2B \\
\midrule
Layers and site & all 12 layers, \texttt{blocks.$L$.hook\_resid\_pre} & layers 5, 12, 19, \texttt{blocks.$L$.hook\_resid\_post} \\
SAEs & in-house ReLU, TopK, JumpReLU (sigmoid surrogate); & public Gemma Scope JumpReLU, width 16{,}384 \\
 & width 3{,}072, target L0 128, 3 seeds (108 SAEs); & (pinned IDs \texttt{average\_l0\_68/82/73}) \\
 & data-scaled R2 SAEs, width 12{,}288; public \texttt{gpt2-small-res-jb} & \\
Datasets & OpenWebText; Geometry of Truth (\texttt{cities}, \texttt{neg\_cities}); RepE-LM & The Pile (validation file) \\
Levels evaluated & reconstruction, information, behaviour, causal & reconstruction, behaviour \\
Strongest evidence & \evlabel{test} (reconstruction), \evlabel{confirmed test} (R2 scaling; R4 behaviour-targeted fine-tuning), & \evlabel{validation} \\
 & \evlabel{validation} (information, causal) & \\
\bottomrule
\end{tabular}
\end{adjustbox}
\end{table}


\paragraph{Models and SAEs.} GPT-2 Small has 12 blocks and a residual stream of width 768. All GPT-2 SAEs read the
activation entering block $L$ (\texttt{blocks.$L$.hook\_resid\_pre}); the beginning-of-sequence (BOS) position is
excluded from training and splicing. Our base SAEs map $x$ to $z=\sigma(W_e(x-b_d)+b_e)$ and $\hat{x}=W_dz+b_d$ with
reconstruction error $e=x-\hat{x}$, using three activation functions: ReLU with an L1 penalty, TopK with $k=128$
\citep{gao2024scaling}, and a JumpReLU whose threshold gradient uses a sigmoid surrogate, which is simpler than the
published estimator \citep{rajamanoharan2024jumping}. Each has 3{,}072 latents and a target L0 (mean active latents per
token) of 128 and is trained at every layer with three seeds, giving 108 SAEs. Architectures are compared only when
every run's L0 is within $\pm10\,\%$ of the target; 107 of 108 runs qualify, and the unmatched layer-0 JumpReLU run is
kept but excluded from rankings. For the scaling experiment (Section~\ref{sec:r2}) we trained SAEs at layers 1, 8 and
11 with 12{,}288 latents on 100M further tokens, together with a width ablation and a random-decoder baseline
\citep{korznikov2026random}. The public \texttt{gpt2-small-res-jb} SAEs \citep{bloom2024gpt2} serve only as an external
reference. In Gemma-2-2B we use the public width-16k Gemma Scope JumpReLU SAEs at layers 5, 12 and 19
\citep{gemmascope2024}, applied to the output of each layer (\texttt{resid\_post}), the site they were trained on, with
model, tokenizer, SAE and data revisions pinned before the run. The two hook conventions follow each SAE family and
imply no layer correspondence between the models. Training details are in Appendix~\ref{app:details}.

\paragraph{Data.} OpenWebText (OWT) documents, split by document, are used to train the base SAEs and to measure
reconstruction and splicing. The Geometry of Truth \texttt{cities} and \texttt{neg\_cities} statements
\citep{marks2023geometry}, split by city, form the truth task; the training split is divided into city groups A and B,
so that training objectives may use only group A and independent evaluation probes are fitted on group B. For the
persona task we built RepE-LM, in which an honest and a dishonest instruction differ in exactly one token at a fixed
position, with honest-synonym and neutral-word controls and held-out wordings. It replaces persona prompts derived from
representation engineering \citep{zou2023representation} whose length alone predicted the label (AUROC 0.80); results
obtained with those prompts are not used. For Gemma we draw 512 evaluation and 64 disjoint reference documents from the
validation file of the uncopyrighted Pile \citep{gao2020pile}, truncated to 1{,}024 tokens. Split sizes, templates and
sampling are in Appendix~\ref{app:details}.

\paragraph{Four levels and their controls.}\label{sec:levels}
\emph{Reconstruction} is measured by explained variance, $\mathrm{EV}=1-\sum\|x-\hat{x}\|^2/\sum\|x-\bar{x}\|^2$, its
complement NMSE, and L0. \emph{Information} is measured by the AUROC gap between a reader applied to original
activations and to reconstructions. A \emph{frozen} reader is fitted on original activations and applied unchanged; a
\emph{re-trained} reader is fitted on reconstructions of training statements. Two references separate SAE-specific
loss from generic perturbation: isotropic noise with the norm of the SAE error, and PCA fitted on OWT activations.
\emph{Behaviour} is measured by splicing $\hat{x}$ into one layer at every non-BOS position and recording the increase
in next-token cross-entropy ($\Delta$CE, nats per prediction), the KL divergence from the original distribution, and
loss recovered relative to mean ablation; a norm-matched random-error control replaces the SAE error by a random
direction of the same norm. Among the four levels, splicing is the only measurement in which the model's forward pass
runs through the SAE output. \emph{Causal preservation} is measured by activation patching between paired prompts;
because the SAE error term is kept, each patch adds a vector to the original activation and does not route
computation through the SAE (Section~\ref{sec:causal}).

\paragraph{Evidence and statistics.}\label{sec:evidence} Each result carries one evidence level: \evlabel{test}
(a held-out split evaluated once under a written protocol), \evlabel{confirmed test} (a pre-registered comparison
first seen on validation and confirmed in one pre-registered held-out evaluation), \evlabel{validation} (a
pre-registered analysis on validation data, which were also used to compare methods, so effect sizes may be
optimistic), \evlabel{pre-registered only} (a fixed protocol that has not produced results) or, where applicable,
\evlabel{exploratory} (not pre-registered, or diagnostic). Unless marked otherwise, results are \evlabel{validation}
findings. Within validation experiments, pre-registered outcome classes such as ``preliminary improvement'', ``no
improvement'', ``not creditable'' (a guard failed) and ``regression'' describe the result of a comparison; they do not
raise or change its evidence level. From the
within-layer reanalysis onwards, hypotheses, criteria and script hashes were written before each run; test splits were
never used for training, tuning or method selection. Seeded GPT-2 comparisons are paired by seed; fine-tuning
comparisons use the pre-registered ``consistent'' rule (better on all three seeds and a mean paired difference above
twice its standard deviation); family comparisons use Welch intervals and require all nine seed pairs to agree.
Gemma intervals are document-bootstrap intervals over one pretrained SAE per layer, not seed variability. Full
procedures and the minimum detectable effects of the design are in Appendix~\ref{app:stats}.
